% Generated draft: requires mathematical and source review.
\documentclass[11pt]{article}
\usepackage[margin=25mm]{geometry}
\usepackage{fontspec}
\setmainfont{Latin Modern Roman}
\usepackage{amsmath,amssymb,mathtools,unicode-math}
\setmathfont{Latin Modern Math}
\usepackage{enumitem,needspace,xurl,hyperref}
\setlength{\parindent}{0pt}
\setlength{\parskip}{5pt}
\setlength{\emergencystretch}{3em}
\newcommand{\N}{\mathbb{N}}
\newcommand{\Z}{\mathbb{Z}}
\newcommand{\Q}{\mathbb{Q}}
\newcommand{\R}{\mathbb{R}}
\newcommand{\C}{\mathbb{C}}
\newcommand{\F}{\mathbb{F}}
\newcommand{\A}{\mathbb{A}}
\newcommand{\PP}{\mathbb{P}}
\newcommand{\E}{\mathbb{E}}
\newcommand{\eps}{\varepsilon}
\newcommand{\dd}{\,\mathrm{d}}
\begin{document}
\section*{UnsolvedMath ID: 700435}
\textbf{Source problem code:} catalogue entry
\subsection{Definitions and mathematical statement}
A profinite group is a topological group isomorphic to an inverse limit of finite discrete groups. A pronilpotent group is an inverse limit of finite nilpotent groups, and a pro-$p$ group is an inverse limit of finite $p$-groups, where $p$ is a prime. Write $\mathfrak N$ for the class of all pronilpotent groups and $\mathfrak P_p$ for the class of all pro-$p$ groups.

Using the standard topological conventions, a \emph{variety of profinite groups} is a nonempty class closed under taking closed subgroups, quotients by closed normal subgroups, and arbitrary Tikhonov products. A subvariety of a variety $\mathfrak U$ is a variety contained in $\mathfrak U$.

For a group $G$, let $\gamma_1(G)=G$ and $\gamma_{n+1}(G)=[\gamma_n(G),G]$, with commutators taken in the underlying abstract group. The group $G$ is \emph{nilpotent} if $\gamma_{c+1}(G)=\{1\}$ for some integer $c\geq 0$. A variety $\mathfrak V$ is \emph{nilpotent}, in the terminology of the supplied source, if every $G\in\mathfrak V$ is nilpotent. No common bound on the nilpotency classes is explicitly included in that definition.

The source calls $\mathfrak V$ \emph{locally nilpotent} if every finitely generated group belonging to $\mathfrak V$ is nilpotent. It does not explicitly specify whether generation is topological or abstract. For the topological reading, put
\[
\mathcal L_{\mathrm{top}}(\mathfrak V)
\quad\Longleftrightarrow\quad
\text{every }G\in\mathfrak V\text{ with }G=\overline{\langle S\rangle}
\text{ for some finite }S\subseteq G\text{ is nilpotent}.
\]
For the literal abstract-generation reading, put
\[
\mathcal L_{\mathrm{abs}}(\mathfrak V)
\quad\Longleftrightarrow\quad
\text{every }G\in\mathfrak V\text{ with }G=\langle S\rangle
\text{ for some finite }S\subseteq G\text{ is nilpotent}.
\]
These readings are retained separately pending clarification of the source's convention.

\textbf{Questions attributed to O. V. Mel'nikov (1980), catalogue entry.} For $\mathcal L$ interpreted according to the intended generation convention, ask:
\begin{enumerate}
\item[(a)] Does every subvariety $\mathfrak V\subseteq\mathfrak N$ containing a non-nilpotent group contain a subvariety $\mathfrak W\subseteq\mathfrak V$ such that $\mathcal L(\mathfrak W)$ holds and $\mathfrak W$ contains a non-nilpotent group?
\item[(b)] For a given prime $p$, does every subvariety $\mathfrak V\subseteq\mathfrak P_p$ containing a non-nilpotent group contain a subvariety $\mathfrak W\subseteq\mathfrak V$ such that $\mathcal L(\mathfrak W)$ holds and $\mathfrak W$ contains a non-nilpotent group?
\end{enumerate}
The supplied record labels the problem partially solved and reports a negative answer to an unspecified subproblem. It does not identify whether that report concerns (a), (b), or both, or specify its dependence on $p$.
\subsection{Short English statement}
Mel'nikov asks whether every non-nilpotent variety of pronilpotent groups contains a locally nilpotent subvariety that is still non-nilpotent, and asks the analogous question for pro-$p$ groups at a fixed prime $p$. The meaning of finite generation needs confirmation. The supplied partial-solution label and negative-answer report do not identify the affected part.
\subsection{Sources}
\begin{itemize}
\item Local UnsolvedMath registry, supplied complete source record: permanent ID $700435$, exact source code \texttt{catalogue entry}, title \emph{Kourovka Notebook Problem 7.38}; upstream revision \texttt{d371ee917ac9647fbb06cd6fc529deec6f3dda82}. Attribution: O. V. Mel'nikov, proposed in 1980.
\item \emph{The Kourovka Notebook (Unsolved Problems in Group Theory)}, 7th issue (1980), Problem 7.38; maintained in the supplied 21st issue edition (2026), arXiv:1401.0300v48. The supplied collection description names E. I. Khukhro and V. D. Mazurov as editors.
\item Supplied discussion and status-review evidence report: a negative answer, described as an AI solution confirmed by E. I. Khukhro, attributed to Achyuth Jayadevan, \emph{Preprint}, 2026, \url{https://zenodo.org/records/22751322}. The supplied evidence does not identify which subproblem this resolves.
\item The supplied status review is dated exactly \texttt{2026-10-16}, assigns canonical status \texttt{partially\_solved}, and describes its evidence as \texttt{maintained\_source}. Its notes attribute the subproblem-solution report to the October 2026 version of the 21st issue. This is supplied dated curation, not an independently verified literature assessment.
\end{itemize}
\end{document}
